\section{Introduction}
\label{sec:introduction}

\openharmony applications use \arkts, a TypeScript-derived language with declarative ArkUI components, framework-managed lifecycle callbacks, and system APIs exposed through both historical \texttt{@ohos.*} and newer \texttt{@kit.*} packages~\cite{openharmony,arktsdoc,harmonyroadmap}. Privacy auditing in this ecosystem requires more than searching for method names. The same member name may denote application code or a platform API; a platform API may be reached through an imported namespace, a manager object returned by a factory, an IR temporary, or a field read rather than a call. Sensitive values may subsequently cross callbacks, Promise handlers, lifecycle entries, logs, storage, or network operations.

ArkAnalyzer provides reusable IR, control-flow, call-graph, and pointer-analysis facilities for ArkTS~\cite{arkanalyzer}. HapFlow builds an IFDS taint analysis on this foundation~\cite{hapflow}. These systems make ArkTS analysis possible, but they leave a distinct front-end problem: a taint solver can propagate only sources that have first been identified. HapFlow's released JSON loader resolves configured namespace and method entries to SDK method signatures. That loader has no representation for a privacy-sensitive property read and does not by itself preserve the package/factory evidence of a call on a manager/helper receiver. Missing that evidence can suppress the initial source fact regardless of the downstream solver.

Consider \texttt{helper.createAsset(...)} after \texttt{helper} is returned by \texttt{photoAccessHelper.getPhotoAccessHelper(ctx)}. The member \texttt{createAsset} alone is insufficient: it is common application vocabulary, while the receiver origin establishes its platform meaning. Conversely, accepting every same-named member in a file importing a media package would create false positives. The analysis must therefore retain the evidence used for a match: namespace-constrained package migration, import alias, receiver origin or type, member, and access kind. A similar problem arises after IR lowering of compound APIs such as \texttt{request.agent.create}, and for fields such as \texttt{deviceInfo.brand}.

We present \mytool, a static analysis system that resolves privacy API identities and assembles provenance-annotated evidence subgraphs for ArkTS. For each candidate usage, \mytool computes a semantic identity over package/namespace migration, receiver evidence, member, and access kind. Indirect calls require an exact SDK type, a compatible call-target signature, a bounded definition/factory-origin chain, or an exact namespace receiver; generic members such as \texttt{on}, \texttt{start}, and \texttt{close} are never accepted from lexical namespace evidence alone. In parallel, the taint query admits only configured sources with an explicit typed carrier and runs IFDS with ArkTS IR recovery; permission-bearing operations without a sensitive output remain detector usages, not taint seeds. A bounded post-IFDS phase supplements selected callback, Promise-handler, and \texttt{async}/\texttt{await} paths; field-stored handlers and ArkUI builder options instead add provenance-labeled report call-graph context. The report keeps detector-local API/sink evidence, privacy-data paths, and framework-input paths distinct, linking detector occurrences only to compatible privacy-data source endpoints.

Our evaluation separates accuracy from prevalence. We first freeze a source-first sample of 60 real projects across three scale strata and two configured-import strata, enumerate executable candidates without reading \mytool reports, and review every candidate at exact source location. The resulting benchmark contains 90 gold occurrences in 16 projects and 44 confirmed zero projects; \mytool recovers all 90 and produces no additional reports. A 64-case positive/negative identity benchmark then isolates the contribution of import, receiver, factory, property, compound-namespace, and package-migration evidence. End-to-end accuracy is measured on all 67 cases in the author-released HapBench suite using its executable artifact as the directly comparable baseline: \mytool attains 100.00\% precision, 94.34\% recall, and 97.09\% F1, with 100\% specificity. A separately specified 48-case asynchronous benchmark attributes nine Promise-\texttt{then} recoveries to post-IFDS supplementation. A single-build 1,014-project study finally measures detector-recognized operations, typed paths, provenance, robustness, and cost without treating corpus output as ground truth.

The evaluation protocol follows lessons from Android analysis research. ReproDroid shows that tools must be compared under the same query and benchmark subset~\cite{reprodroid}; TaintBench stores expected and unexpected flows rather than treating every undiscovered report as a false positive~\cite{taintbench}; and mutation studies show that optimizations can introduce undocumented unsoundness~\cite{muse}. Accordingly, we publish all HapBench labels and errors, distinguish incomplete runs from zero-flow results, and report precision, recall, specificity, balanced accuracy, F1, MCC, and per-category behavior. Android analyzers such as FlowDroid and PacDroid inform our design, but are not presented as numerical baselines because they analyze Android bytecode and Android framework models rather than ArkTS~\cite{arzt14,pacdroid}.

This paper makes the following contributions:
\begin{itemize}
    \item \textbf{Evidence-carrying API identity resolution.} We formulate privacy API recognition as a package-, receiver-, member-, and access-aware matching problem, covering direct and assigned calls, manager/helper receivers, IR aliases, and executable property reads while recording the acceptance witness of each match.
    \item \textbf{Carrier-aware ArkTS propagation.} We define typed source carriers and seed facts for privacy-data returns and callbacks, distinguish framework-provided inputs, and reject detector-only operations as taint seeds. ArkTS IR recovery operates inside IFDS, while bounded callback/Promise rules supplement unresolved asynchronous paths after the solver.
    \item \textbf{Provenance-annotated audit subgraphs.} We retain API identity, call context, detector-local sinks, configured-query paths, permissions, and source locations in one artifact while preserving which subsystem produced each edge and whether a configured source was linked to a detector usage.
    \item \textbf{A reproducible multi-level evaluation.} We contribute a source-first real-project occurrence benchmark, a matched 64-case identity benchmark, a complete reproduction on the official 67-case HapBench suite, a balanced 48-case asynchronous benchmark, and a clone-aware study of 1,014 completed projects. Each layer has a distinct oracle and claim boundary, while machine-readable manifests connect reported values to the evaluated build, rules, SDK, and project inventory.
\end{itemize}
